% Figure 1, panel D: the graph prior, main-paper eq (graph_prior). A small graph
% with its edges coloured by the evidence class of the pair they occupy, and the
% prior written out beneath it. Every ordered pair is an edge of D independently
% with the edge probability of its own class, so what the picture has to show is that
% the class travels with the pair, not with the node.
%
% The graph is deliberately untyped: no TFs, enhancers or genes, and no node
% labels. The graph prior does not depend on what the nodes are, only on which
% class each pair falls in, and typing the nodes would import a constraint the
% prior does not have. Panel C makes the same choice for the same reason. The
% layout nonetheless survives a typed reading -- two nodes per column, and the
% classes on the edges fix the columns as TFs, enhancers and genes -- so a reader
% who carries panel B's sources over to this panel finds every drawn pair
% legitimate rather than having to un-type the drawing first.
%
% Two permitted pairs are drawn dashed and absent. Without them the panel shows
% only the e_ij = 1 half of the prior, and the (1 - rho) factors in the equation
% would have nothing in the picture answering to them. One is baseline and one
% carries evidence: with only the baseline pair dashed the picture would say that
% a pair a source supports is an edge, which is exactly what SPARK does not
% assume. rho_k is a probability and not a rule, and calibrating it is the point.
%
% Source 4 carries two edges rather than one. With a single edge per class the
% picture reads as though a class were a pair, and the one thing the graph prior
% has to say is the opposite: the class travels with the pair, and every pair in a
% class shares that class's one probability. Two edges in the same colour say it, and the
% rho_4 squared in the equation is what it costs there.
%
% The two notes under the graph are drawn rather than left to the caption, since
% they are what makes the equation readable off the picture: what the dashed pairs
% mean, and that every pair not drawn is forbidden. The second is not decoration:
% rho_{-1} = 0, so a forbidden pair contributes (1-0)^1 = 1 and drops out, which is
% what makes the eight factors below the whole product rather than a selection.
%
% Deliberately NOT here: the conjugate conditional rho_k | D. That is a
% consequence of the model rather than part of its specification, and it is the
% Gibbs step of the sampler, so it sits in the MCMC band between the model row
% and the output. Panel E carries the Beta priors on the class edge probabilities.
%
% Colour means evidence class everywhere in this figure, so the drawn edges take
% one colour each, baseline included.
% Local origin is the panel's top-left corner. Panel size 5.25 x 4.35.
\draw[panel] (0,0) rectangle (5.25,-4.35);
\node[ttl] at (0.15,-0.12) {D\quad Graph prior};

%%%%%%%%%%%%%%%%%%%%%%%%%%%%%%%%%%%%%%%%%%%%%%%%%%%%%%%%%%%%%%%%%%%%%%%%%%%%%%%%
% A small graph, its edges coloured by the evidence class of the pair
%%%%%%%%%%%%%%%%%%%%%%%%%%%%%%%%%%%%%%%%%%%%%%%%%%%%%%%%%%%%%%%%%%%%%%%%%%%%%%%%
\begin{scope}[shift={(0,0)}]
  \coordinate (ga) at (0.85,-1.18);   % supported by sources 1 and 2
  \coordinate (gb) at (0.85,-2.38);
  \coordinate (gc) at (2.62,-1.33);   % the middle node the sources support
  \coordinate (gd) at (2.62,-2.23);   % reached only by baseline pairs
  \coordinate (ge) at (4.40,-1.18);
  \coordinate (gf) at (4.40,-2.38);

  % The two permitted pairs that are not edges of D, each drawn in its own class's
  % colour rather than a neutral grey: that is what makes the (1 - rho) factors in
  % the equation below readable off the picture. Arrowheaded like the edges, since
  % each stands for one ordered pair and not for the two between that pair of
  % nodes. Without the head a reader may read one as both, and the equation would
  % owe them (1 - rho) squared rather than once.
  \draw[->, clsBase, dashed, dash pattern=on 1.4pt off 1.2pt, line width=0.9pt,
        shorten >=4.2pt, shorten <=4.2pt] (gb) -- (gd);
  \draw[->, clsSrc3, dashed, dash pattern=on 1.4pt off 1.2pt, line width=0.9pt,
        shorten >=4.2pt, shorten <=4.2pt] (gd) -- (ge);

  % The edges of D. The two baseline pairs sit on the same second middle node, one
  % an edge and one not: the baseline class is the permitted pairs that no source
  % supports, and putting both on one node says that rho_0 is what decides between
  % them.
  \draw[->, clsBase, line width=1.1pt, shorten >=4.2pt, shorten <=4.2pt] (ga) -- (gd);
  \draw[->, clsSrc1, line width=1.1pt, shorten >=4.2pt, shorten <=4.2pt] (ga) -- (gc);
  \draw[->, clsSrc2, line width=1.1pt, shorten >=4.2pt, shorten <=4.2pt] (gb) -- (gc);
  \draw[->, clsSrc3, line width=1.1pt, shorten >=4.2pt, shorten <=4.2pt] (gc) -- (ge);
  \draw[->, clsSrc4, line width=1.1pt, shorten >=4.2pt, shorten <=4.2pt] (gc) -- (gf);
  \draw[->, clsSrc4, line width=1.1pt, shorten >=4.2pt, shorten <=4.2pt] (gd) -- (gf);

  \foreach \n in {ga,gb,gc,gd,ge,gf}{
    \draw[fill=white, draw=black!45, line width=0.6pt] (\n) circle (0.135);
  }

  % Clear of the lower nodes, which sit at -2.38 with radius 0.135.
  \node[note, text=black!55] at (2.625,-2.78) {dashed: permitted, not an edge};
  \node[note, text=black!55] at (2.625,-3.02) {pairs not drawn: forbidden};
\end{scope}

%%%%%%%%%%%%%%%%%%%%%%%%%%%%%%%%%%%%%%%%%%%%%%%%%%%%%%%%%%%%%%%%%%%%%%%%%%%%%%%%
% The prior, evaluated for the graph above rather than stated in general
%%%%%%%%%%%%%%%%%%%%%%%%%%%%%%%%%%%%%%%%%%%%%%%%%%%%%%%%%%%%%%%%%%%%%%%%%%%%%%%%
% Each rho is set in its own class colour, so the factor a reader is looking at
% names the edge it came from without the panel having to label the edges too.
% The baseline is mixed toward black first: at #BFC6D0 it reads as a line but
% not as type.
% Set on two lines: with the second dashed pair and the squared factor the product
% runs about 5.5 cm at \scriptsize, wider than the 5.25 cm panel. The break is
% after the edges, so the first line is the pairs that are edges and the second
% the pairs that are not.
\node[anchor=north, inner sep=0pt] at (2.625,-3.26) {\scriptsize $
  \begin{aligned}
    P(\mathcal{D} \mid \boldsymbol{\rho}) &\propto
      \textcolor{clsBase!55!black}{\rho_0}\,
      \textcolor{clsSrc1}{\rho_1}\,
      \textcolor{clsSrc2}{\rho_2}\,
      \textcolor{clsSrc3}{\rho_3}\,
      \textcolor{clsSrc4}{\rho_4^2}\\[2pt]
    &\quad\times \bigl(1 - \textcolor{clsBase!55!black}{\rho_0}\bigr)
      \bigl(1 - \textcolor{clsSrc3}{\rho_3}\bigr)
  \end{aligned}$};
